% ADDED figure (2026-11-02, Lecture 16 merge); generated by qft-C/tools/l16/figs/make_figs.sh
\tikzset{
  fdl/.style={line width=0.8pt, black},
  fdx/.style={circle, fill=black, inner sep=0pt, minimum size=3.6pt},
  fdv/.style={circle, fill=figB, draw=figB, inner sep=0pt, minimum size=5pt},
  fdlab/.style={font=\scriptsize, inner sep=2.6pt},
  fdmom/.style={->, >=Stealth, figO, line width=0.6pt},
}
\tikzset{
  fda/.style={fdl, postaction={decorate}, decoration={markings, mark=at position 0.55 with {\arrow{Stealth[length=5pt]}}}},
  fdS/.style={font=\small},
}
\begin{tikzpicture}[font=\small]
  \tikzset{dl/.style={line width=0.8pt, dashed, figR}}
  % (1) zeroth order: delta lines
  \begin{scope}
    \draw[dl] (-0.9,0.55)--(0.9,0.55); \draw[dl] (-0.9,-0.55)--(0.9,-0.55);
    \node[fdlab,left] at (-0.9,0.55){$A$}; \node[fdlab,left] at (-0.9,-0.55){$B$};
    \node[fdlab,right] at (0.9,0.55){$1$}; \node[fdlab,right] at (0.9,-0.55){$2$};
    \node[fdS,align=center] at (0,-1.55) {(1) $\lambda^0$: lines are $\delta$'s\\[1pt] part of the $\mathbf 1$};
  \end{scope}
  % (2) bubble times delta lines
  \begin{scope}[xshift=4.4cm]
    \draw[fdl] (-1.15,0.3) circle (0.3); \draw[fdl] (-1.15,-0.3) circle (0.3); \node[fdv] at (-1.15,0){};
    \node at (-0.6,0) {$\times$};
    \draw[dl] (-0.2,0.55)--(1.2,0.55); \draw[dl] (-0.2,-0.55)--(1.2,-0.55);
    \node[fdlab,left] at (-0.2,0.55){$A$}; \node[fdlab,left] at (-0.2,-0.55){$B$};
    \node[fdlab,right] at (1.2,0.55){$1$}; \node[fdlab,right] at (1.2,-0.55){$2$};
    \node[fdS,align=center] at (0,-1.55) {(2) $\lambda$: bubble $\times$ $\delta$'s\\[1pt] bubble omitted; the $\mathbf 1$};
  \end{scope}
  % (3) tadpole on a passing line
  \begin{scope}[xshift=8.8cm]
    \draw[fdl] (-0.9,0.55)--(0.9,0.55); \draw[fdl] (0,0.85) circle (0.3); \node[fdv] at (0,0.55){};
    \draw[dl] (-0.9,-0.55)--(0.9,-0.55);
    \node[fdlab,left] at (-0.9,0.55){$A$}; \node[fdlab,left] at (-0.9,-0.55){$B$};
    \node[fdlab,right] at (0.9,0.55){$1$}; \node[fdlab,right] at (0.9,-0.55){$2$};
    \node[fdS,align=center] at (0,-1.55) {(3) $\lambda$: two pieces, $p_1 = p_A$\\[1pt] forward: part of the $\mathbf 1$};
  \end{scope}
  % (4) correction on an external leg (slide 26)
  \begin{scope}[xshift=13.2cm]
    \draw[fdl] (-0.95,0.75)--(-0.45,0.4); \draw[fdl] (-0.45,0.4)--(0,0); \draw[fdl] (-0.31,0.57) circle (0.22);
    \node[fdv] at (-0.45,0.4){};
    \draw[fdl] (-0.95,-0.75)--(0,0); \draw[fdl] (0,0)--(0.95,0.75); \draw[fdl] (0,0)--(0.95,-0.75); \node[fdv] at (0,0){};
    \node[fdlab,left] at (-0.95,0.75){$A$}; \node[fdlab,left] at (-0.95,-0.75){$B$};
    \node[fdlab,right] at (0.95,0.75){$1$}; \node[fdlab,right] at (0.95,-0.75){$2$};
    \node[fdS,align=center] at (0,-1.55) {(4) $\lambda^2$: loop on leg $A$\\[1pt] state correction: amputated};
  \end{scope}
  % (5) connected amputated
  \begin{scope}[xshift=17.4cm]
    \draw[fdl] (-0.8,0.7)--(0,0)--(0.8,0.7); \draw[fdl] (-0.8,-0.7)--(0,0)--(0.8,-0.7); \node[fdv] at (0,0){};
    \node[fdlab,left] at (-0.8,0.7){$A$}; \node[fdlab,left] at (-0.8,-0.7){$B$};
    \node[fdlab,right] at (0.8,0.7){$1$}; \node[fdlab,right] at (0.8,-0.7){$2$};
    \node[fdS,align=center] at (0,-1.55) {(5) $\lambda$: connected, amputated\\[1pt] $i\hat T$: $i\mathcal M = -i\lambda$};
  \end{scope}
\end{tikzpicture}
